% ============================================================
\chapter{Conclusion}
\label{chap:conclusion}
% ============================================================

\section{Conclusion}

PipedPeer runs an unmodified Python script on other machines with no cluster to
configure, reproduces its environment exactly, distributes the parallel work
inside it without the user touching the source, and joins machines on different
networks with a single command and no relay. Its six contributions are:

\begin{enumerate}[leftmargin=*, label=\arabic*.]
  \item \textbf{Peer-to-peer closure distribution.} Pinning an exact package set
  makes an environment's store path identical on every machine and on any day,
  so a worker that holds an environment pays nothing for it (0.83\,s warm
  against 23.4\,s cold), one that holds part of it receives only the paths it
  lacks (86\% fewer bytes), and signed closures move between peers.
  \item \textbf{Transparent interception with measured work sharing.} The
  parallel primitives a script already uses are substituted at import time;
  parallel maps are divided by a makespan solver over measured peer rates; and
  numerical work moves only when the cost model says moving it is cheaper than
  computing it.
  \item \textbf{Joining across networks without a relay.} An introducer in every
  daemon that never carries data, router port mapping, a candidate race over one
  shared socket, a birthday-collision fallback for the hardest NATs, and
  authenticated QUIC links that the rest of the system sees as ordinary nodes.
  \item \textbf{Placement without a scheduler.} Filters and a bounded score are
  evaluated in the submitting client, under a three-phase lease whose expiry is
  the crash recovery.
  \item \textbf{A fallback that preserves the result on every distributed path.}
  Every remote path has a working local path beneath it, so a script with no
  failure handling never sees a failure.
  \item \textbf{Speed-weighted data-parallel training through the daemon.}
  Gradients are averaged by samples in 8-bit with error feedback, batches are
  sized by measured speed, and exact or fast synchronisation is chosen by ring
  size; at 32 steps on two ranks, fast synchronisation held accuracy while each
  rank sent 7.6 to 9.2 times fewer bytes.
\end{enumerate}

The principle that runs through the design is that the cluster is used only when
a measurement says it will help: work is divided by measured speed rather than
by headcount, dense matrix products stay local on a CPU cluster, a training ring
whose size the user did not fix is released when one GPU is faster, and fast
synchronisation is chosen only for ring sizes where it kept accuracy.

\section{Limitations}

\begin{itemize}[leftmargin=*]
  \item \textbf{No relay.} Two machines between which no direct path exists ---
  for example two carrier networks that both allocate a new port per
  destination --- cannot work together; such a peer is reported and routed
  around.
  \item \textbf{Python and Linux only}, with NVIDIA GPUs in practice; phones
  take parallel chunks only.
  \item \textbf{Submitter and lead.} A parallel script runs on its submitter,
  which must survive the job, and a training ring needs rank 0's machine.
  \item \textbf{Fast synchronisation for small rings.} No interval kept accuracy
  at eight ranks, so automatic mode keeps rings of five or more exact.
  \item \textbf{Interception gaps.} Only \texttt{map} of a process-pool executor
  is intercepted, the \texttt{@} operator is not, and numerical work under
  32\,MiB never leaves the node, by design.
  \item \textbf{A partial sandbox.} The job runs as user 0 inside the container
  with the host's network and no seccomp filter, and isolation is skipped with a
  warning on a node lacking \texttt{crun} or user namespaces.
  \item \textbf{Measurement.} Speedup comes from containers with fixed processor
  quotas and a few real machines; the success rate of hole punching across many networks, two
  GPUs in one machine and fast synchronisation at five to seven ranks are not
  measured. NAT traversal is evaluated on the two-site testbed on one host and
  over an overlay between real homes.
  \item \textbf{Seed spread.} The 1\% accuracy bar applies to the mean over
  seeds; single seeds spread wider, from $-3.9\%$ to $+2.2\%$ against one
  device, so one run can miss the bar while the mean meets it.
  \item \textbf{Trust.} Every holder of the cluster secret can run code on every
  node (Section~\ref{sec:security}).
\end{itemize}
